\scanpage{48}
Soit
\[
 Q=\operatorname{Var}(\unsure{\Delta|\phi})\subset\mathbb P^\vee(\mathcal M)\tag*{(4)}
\]
où $\Delta$ est la forme quadratique
\[
 \Delta(u)=-\det u\tag*{(5)}
\]
sur $\mathcal M$ -- qui est non dégénérée. On va définir un isomorphisme canonique
\[
 Q\xrightarrow{\overset{\psi}{\sim}}\deleted{X\times X}\ \mathbb P(V)\times\mathbb P(V)\quad(\simeq X\times X)\tag*{(6)}
\]
En effet, $Q$ correspond aux \deleted{\unsure{$u:\in\Gamma(\mathcal M)$|$u\in\Gamma(\mathcal M)$} partout} \deleted{\unsure{$V\longrightarrow V$}}
\[
 u:V\mathrel{\unsure{-|\longrightarrow}}V
\]
(non nuls sur chaque fibre) tels que $u$ soit \uline{\unsure{strictemt|strictement} de rang 1} i.e. se factorise canoniquement en
\diagram{\begin{tikzcd}[ampersand replacement=\&] V\arrow[r,"\text{épi}"]\&I\arrow[r,"\substack{\text{mono}\\\text{univ}}"]\&V\end{tikzcd}}
avec $I$ inversible, et il est connu à \inserted{multipli\unsure{c.|.}} scalaire près quand on connaît le quotient $I$ de $V$, \deleted{\unsure{et le ss.}} ou encore le noyau de $u$, et \deleted{l'image} le sous-module $I$ de $V$ i.e $\operatorname{Im}u$. Ceux-ci peuvent être pris \unsure{indépendam$^{\mathrm t}$|indépendants|indépendamment}. Donc (6) est induit par
\[
 \underset{\substack{\deleted{\unsure{m}}\\\text{non nul sur}\\\text{tte fibre,}\\\Delta(u)=0}}{u}
 \longmapsto\bigl(\underset{\substack{\shortparallel\\L}}{\operatorname{Ker}u},
 \underset{\substack{\shortparallel\\M}}{\operatorname{Im}u}\bigr).\tag*{(7)}
\]
% slanted
\margin{\uline{NB} L'isom en sens inverse est encore plus évident, savoir
$(L,M)\mapsto L^\perp\otimes M$\\
$\hfill\rotatebox{90}{$\subset$}\qquad$\\
$V^\vee\otimes V\simeq\mathcal M$}
% archive sheet number: 47
